% /root/monodromy, 19 November 2026, 05:39 UTC.
\subsection{An ambient adjunction obstruction in characteristic two}
\label{subsec:general-adjunction-covector}
\status{Proved under the stated hypotheses; independently audited.}

\begin{theorem}\label{thm:general-adjunction-covector}
Let $X/k$ be smooth, projective and geometrically connected, of
dimension $n\geq2$, over a perfect field of characteristic two.
Let $L=\mathcal O_X(1)$ for the fixed embedding $X\subset\mathbf P^r$.
Write $R_I=\operatorname{reg}(\mathcal I_X)$.
Suppose that, for some
integer $d_*$, there is a line bundle $M_*$ defined over $k$ with
\[
 M_*^{\otimes2}\simeq\omega_X\otimes L^{d_*}.
\]
Put $d=d_*+2t$ and $M=M_*\otimes L^t$.
Write $F=F_{X/k}:X\to X^{(2)}$ and
$\mathcal B=\operatorname{coker}(\mathcal O_{X^{(2)}}\to F_*\mathcal O_X)$.
Choose Castelnuovo--Mumford regularities
\[
 R_M=\operatorname{reg}_L(M_*),\qquad
 R_B=\operatorname{reg}_{L^{(2)}}
       (\mathcal B^\vee\otimes M_*^{(2)}),
 \qquad
 \kappa=\frac{c_1(\omega_X)c_1(L)^{n-1}}{c_1(L)^n}.
\]
Assume
\begin{equation}\label{eq:adjunction-covector-threshold}
 t\geq\max\{0,R_M,R_B-1\},\qquad
 d\geq\max\{d_{\rm coh},R_I-1\},\qquad d>\kappa.
\end{equation}
On the maximal-stable-rank smooth-hypersurface open $U_d$,
the relative Hasse--Witt local system $\mathcal V_d$ has a
canonical nonzero morphism
\begin{equation}\label{eq:adjunction-invariant-covector}
 \lambda_M:\mathcal V_d\longrightarrow\mathbf F_2.
\end{equation}
Here ``canonical'' is relative to the given square root and its
isomorphism with $\omega_X L^d$; the construction is described below.
In particular, both the geometric and arithmetic images fix a
nonzero covector.  If $\operatorname{rank}\mathcal V_d\geq2$,
they are proper subgroups of $\GL(\mathcal V_{d,\bar\eta})$ and
the representation is reducible.

If the square root is defined only after an algebraic extension
of $k$, the conclusion about the geometric image still holds.
There are arbitrarily large degrees satisfying
\eqref{eq:adjunction-covector-threshold}; the rank is at least two
for all sufficiently large such degrees, by the previously proved
explicit stable-rank lower bound (or by the all-residue theorem
when $k$ is finite).
\end{theorem}

\begin{proof}
The morphism $F$ is finite flat of degree $2^n$.  Finite duality,
or the Cartier trace calculation of
\cite[Lemma~1.3.6, Proposition~1.3.7, Theorem~1.3.8, pp.~26--28]
{BrionKumar2005}, gives
\begin{align}
 F^!M^{(2)}
 &\simeq \omega_X\otimes F^*(\omega_{X^{(2)}}^{-1}\otimes M^{(2)})
       \simeq\omega_X^{-1}\otimes M^{\otimes2}
       \simeq L^d,\label{eq:adjunction-finite-duality}\\
 H^0(X,L^d)
 &\simeq\Hom_{X^{(2)}}(F_*\mathcal O_X,M^{(2)}).
 \label{eq:adjunction-section-map}
\end{align}
These are relative-Frobenius identifications and are $k$-linear.
Given $s\in H^0(X,L^d)$, let $\tau_s$ be the corresponding map.
Set
\begin{equation}\label{eq:adjunction-w-definition}
 v_s=\tau_s(1)\in H^0(X^{(2)},M^{(2)}),\qquad
 w_s=F^*v_s\in H^0(X,M^{\otimes2})
                  =H^0(X,\omega_X L^d).
\end{equation}
The map $s\mapsto w_s$ is $k$-linear.  Relative Frobenius is
essential in this assertion: its pullback fixes base-field
coefficients.  On a perfect geometric fibre one can express
$w_s$ as a square using the absolute Cartier operator, but
one must not square the universal parameter coefficients.

We check the Frobenius identity on geometric fibres.  Locally
choose a generator $m$ of $M$ and etale coordinates whose top
form is $\eta$.  The induced generator of $L^d$ is $m^2/\eta$;
write $s=fm^2/\eta$.  Write $T$ for the coefficient form of
the absolute Cartier trace, so that
\[
 C(f\eta)=T(f)\eta,\qquad T(a^2b)=aT(b).
\]
Formula~\eqref{eq:adjunction-w-definition} becomes
$w_s=T(f)^2m^2$.  The Cartier operator dual to
\[
 F_s:H^n(X,L^{-d})\longrightarrow H^n(X,L^{-d}),
       \quad \xi\longmapsto sF_X(\xi),
\]
sends $gm^2$ to $T(fg)m^2$.  Consequently
\begin{equation}\label{eq:adjunction-covector-Cartier}
 C_s(w_s)=T(fT(f)^2)m^2=T(f)^2m^2=w_s.
\end{equation}
Under Serre duality this is exactly
\[
 \langle F_s\xi,w_s\rangle=\langle\xi,w_s\rangle^2.
\]
Thus $w_s$ is a Frobenius-equivariant scalar-valued covector.
In matrix notation this is $A^tw_s=w_s^{[2]}$, consistently
with the explicit projective-space formula.

Because $w_s$ is linear in $s$, it defines a global covector
on $H^n(X,L^{-d})\otimes\mathcal O(-1)$ over
$\mathbf P H^0(X,L^d)$.  It restricts to the relative coherent
bundle
\[
 \mathcal E_d=\ker\bigl(H^n(X,L^{-d})\otimes\mathcal O(-1)
       \xrightarrow{\ s\ }H^n(X,\mathcal O_X)\otimes\mathcal O\bigr).
\]
The Frobenius identity shows that the restricted covector
kills the Frobenius-nilpotent part and induces a map from
the stable Lang local system to the constant sheaf
$\mathbf F_2$.

It remains to prove nonvanishing.  The sheaf $\mathcal B$ is
locally free: etale locally the standard square-free
monomials form a basis of $F_*\mathcal O_X$ containing $1$.
Dualizing its defining exact sequence and tensoring with
$M^{(2)}$ gives
\[
 0\longrightarrow\mathcal B^\vee\otimes M^{(2)}
 \longrightarrow\mathcal Hom(F_*\mathcal O_X,M^{(2)})
 \xrightarrow{\ \tau\mapsto\tau(1)\ } M^{(2)}
 \longrightarrow0.
\]
The regularity bounds in
\eqref{eq:adjunction-covector-threshold} give
\[
 H^1(X^{(2)},\mathcal B^\vee\otimes M^{(2)})=0,
       \qquad H^0(X^{(2)},M^{(2)})\ne0.
\]
Indeed, the first is the degree-one regularity vanishing,
and the second follows from global generation of the
nonzero line bundle $M$.  Thus $s\mapsto v_s$ is surjective
onto a nonzero space, and $w_s\ne0$ for generic $s$.

Take such an $s$ with $Y=(s=0)$ smooth.  If its covector
vanished on $E_s$, duality for the displayed kernel would
give $w_s=s\eta_0$ for some $\eta_0\in H^0(X,\omega_X)$.
Every valuation of $w_s=F^*v_s$ is even.  The divisor of
$s$ is reduced, so the valuation of $\eta_0$ along each
component of $Y$ would be odd and positive.  Hence
$\eta_0\in H^0(X,\omega_X L^{-d})$.  This last space is zero:
a nonzero section would give an effective divisor with
intersection number $(\kappa-d)c_1(L)^n<0$ against
$L^{n-1}$, an impossibility.  Thus the restricted covector
is nonzero generically.  A morphism of finite local
systems on the connected $U_d$ has locally constant rank;
it is therefore nonzero everywhere on the maximal-rank open
of the complete linear system.

For the prescribed ambient degree-$d$ hypersurface parameter
space, the additional bound $d\geq R_I-1$ gives
$H^1(\mathbf P^r,\mathcal I_X(d))=0$.  The restriction
$H^0(\mathbf P^r,\mathcal O(d))\to H^0(X,L^d)$ is therefore
surjective.  Its induced linear projection, away from its
center, pulls back the construction just given.  The smooth
ambient open avoids the center; its maximal stable rank is
the same as in the complete linear system, and the pulled-back
covector is nonzero on it.  This is the $U_d$ in the statement.

Everything was defined over $k$ if the square root was,
so arithmetic as well as geometric monodromy fixes it.
After extending the base field the same construction
gives the asserted geometric conclusion.  The kernel
of a nonzero covector is a nonzero proper invariant
subspace when the rank is at least two.
\end{proof}

\begin{remark}\label{rem:adjunction-square-Pn}
For $X=\mathbf P^n$ and $L=\mathcal O(1)$, the square-root
condition is $d-n-1\equiv0\pmod2$.  In homogeneous
monomial coordinates the construction recovers the
covector whose entry indexed by $u$ is the defining
coefficient $c_u$ when all coordinates of $u$ are odd,
and zero otherwise.  Thus this obstruction is an
ambient adjunction phenomenon, rather than a property
limited to plane theta characteristics.
The theorem gives a subgroup constraint for general
$X$, not equality with the full covector stabilizer.
\end{remark}
